% glossary.tex -- Module 6 glossary (A-Z). GENERATED by
\unnumberedchapter[\faBook]{Glossary}

Every technical word in this book, from A to Z. The meaning is in plain English. The
last column gives the usual Chinese term used in physics textbooks.

\begin{glossarytable}
  % ---- A
  \gterm{AC induction motor}{a motor whose rotor currents are induced by a rotating field, not fed in by wires}{交流感应电动机（异步电动机）}
  \gterm{Acceleration}{how fast the velocity changes}{加速度}
  \gterm{Accuracy}{how close a result is to the true value}{准确度}
  \gterm{Action--reaction pair}{two equal and opposite forces that two objects put on each other (Newton's third law)}{作用力与反作用力}
  \gterm{Alternating current (AC)}{current that keeps reversing its direction}{交流电}
  \gterm{Ampere (A)}{the SI unit of current: \qty{1}{A} means one coulomb of charge passes a point each second}{安培}
  \gterm{Attract / repel}{pull together / push apart}{吸引／排斥}
  \gterm{Average emf}{the emf found from the whole change in flux and the whole time it takes}{平均感应电动势}
  \gterm{Axle}{the rod through the middle of a coil that it turns on}{转轴}
  % ---- B
  \gterm{Back emf ($\varepsilon_\mathrm{back}$)}{the emf induced in a spinning motor coil; it opposes the supply}{反电动势}
  \gterm{Beam}{a narrow stream of fast particles, such as electrons from an electron gun}{粒子束（电子束）}
  \gterm{Bearings}{the parts that hold a spinning axle and let it turn smoothly}{轴承}
  \gterm{Brushes}{fixed contacts that join the spinning coil to the circuit}{电刷}
  % ---- C
  \gterm{Cathode}{the negative electrode; when it is heated, it gives off electrons}{阴极}
  \gterm{Centripetal force}{the net force towards the centre of a circle}{向心力}
  \gterm{Charge-to-mass ratio ($q/m$)}{charge divided by mass; a bigger $q/m$ gives a bigger acceleration}{荷质比（比荷）}
  \gterm{Coil}{many loops (turns) of wire wound together}{线圈}
  \gterm{Conventional current}{the way positive charge would flow: from $+$ round the circuit to $-$}{（规定的）电流方向}
  \gterm{Core (soft-iron core)}{the iron that carries the magnetic field from one coil to the other}{铁芯（软铁芯）}
  \gterm{Couple}{two equal and opposite forces that turn an object without moving it sideways}{力偶}
  \gterm{Crossed fields}{an electric field and a magnetic field at right angles in the same place}{正交电磁场}
  \gterm{Current loop}{one turn of wire carrying a current; it acts like a flat magnet}{环形电流}
  \gterm{Cycle}{one full swing of AC, forwards and back}{周（一次循环）}
  % ---- D
  \gterm{DC motor}{a device that turns electrical energy from a DC supply into rotational kinetic energy}{直流电动机}
  \gterm{Deflection}{how far a particle is pushed sideways from its straight path}{偏转（偏移量）}
  \gterm{Direct current (DC)}{current that flows in one direction only}{直流电}
  \gterm{Directly proportional ($\propto$)}{double one and the other doubles; the graph is a straight line through the origin}{成正比}
  % ---- E
  \gterm{Eddy currents}{loops of current induced inside solid metal when the flux through it changes}{涡流}
  \gterm{Edge effect}{the slight bulging of field lines near the ends of charged plates}{边缘效应}
  \gterm{Effective value (of AC)}{the steady (DC) voltage that would heat a heater by the same amount}{有效值}
  \gterm{Efficiency ($\eta$)}{the percentage of the input power that comes out as useful power}{效率}
  \gterm{Electric field}{a region where a charge feels an electric force}{电场}
  \gterm{Electric field strength ($E$)}{the force on each coulomb of charge; unit \si{V.m^{-1}}, the same as \si{N.C^{-1}}}{电场强度}
  \gterm{Electric potential energy ($U$)}{energy a charge stores because of its position in an electric field}{电势能}
  \gterm{Electricity grid}{the network of power stations, transmission lines and transformers}{电网}
  \gterm{Electromagnetic induction}{making an emf by changing the magnetic flux through a circuit}{电磁感应}
  \gterm{Electromotive force (emf, $\varepsilon$)}{not a force: the energy a source gives each coulomb of charge; unit volt (V)}{电动势}
  \gterm{Electron flow}{the real movement of electrons in a wire, opposite to conventional current}{电子的定向移动方向}
  \gterm{Electron gun}{plates with a voltage that speed electrons up from rest into a beam}{电子枪}
  \gterm{Electron volt (eV)}{the energy one electron gains through \qty{1}{V}: \qty{1}{eV} $=$ \qty{1.602e-19}{J}}{电子伏特}
  \gterm{Electronic balance}{a digital scale; it shows the force on its pan as a mass}{电子天平}
  \gterm{Elementary charge ($e$)}{the size of the charge on one proton or electron, \qty{1.602e-19}{C}}{元电荷}
  \gterm{External magnetic field}{a field made by something else (such as a magnet), not by the wire itself}{外磁场}
  % ---- F
  \gterm{Faraday's law}{the induced emf equals the number of turns times the rate of change of flux}{法拉第电磁感应定律}
  \gterm{Ferrite}{a magnetic ceramic that conducts electricity very poorly}{铁氧体}
  \gterm{Field lines}{lines that show the direction of a field; closer lines mean a stronger field}{电场线；磁感线}
  \gterm{Force per unit length ($F/l$)}{the force on each metre of wire; unit \si{N.m^{-1}}}{单位长度导线受的力}
  \gterm{Frequency ($f$)}{the number of full cycles each second; unit hertz (Hz)}{频率}
  % ---- G
  \gterm{Galvanometer}{a sensitive meter whose needle can move either way from zero}{灵敏电流计}
  \gterm{Generator}{a device that turns kinetic energy into electrical energy}{发电机}
  \gterm{Gradient}{how steep a graph is: rise divided by run}{斜率}
  % ---- H
  \gterm{Helix}{a spiral path, like a stretched spring}{螺旋线}
  \gterm{Hypotenuse}{the longest side of a right-angled triangle}{斜边}
  \gterm{Hysteresis loss}{energy turned into heat each cycle as the core is magnetised one way, then the other}{磁滞损耗}
  % ---- I
  \gterm{Ideal transformer}{a perfect transformer that wastes no energy, so power in equals power out}{理想变压器}
  \gterm{In series}{joined one after another in one loop, so the same current flows through each part}{串联}
  \gterm{Incomplete flux linkage (flux leakage)}{some of the flux from the primary misses the secondary}{漏磁（磁通未完全耦合）}
  \gterm{Induce}{to cause something without touching it; a changing flux induces an emf}{感应（产生）}
  \gterm{Induced current}{the current that an induced emf drives around a closed circuit}{感应电流}
  \gterm{Induced emf}{an emf produced by a changing magnetic flux}{感应电动势}
  \gterm{Induced magnetic field}{the magnetic field made by an induced current}{感应电流的磁场}
  \gterm{Ion}{an atom with extra or missing electrons}{离子}
  % ---- L
  \gterm{Laminated core}{a core made of thin iron sheets with insulation between them}{叠片铁芯（硅钢片铁芯）}
  \gterm{Law of conservation of energy}{energy cannot be created or destroyed; it only changes form}{能量守恒定律}
  \gterm{Lenz's law}{an induced current's magnetic field opposes the change in flux that made it}{楞次定律}
  \gterm{Lever arm}{the perpendicular distance from the axis to the line of a force}{力臂}
  \gterm{Like currents}{currents in the same direction; parallel like currents attract}{同向电流}
  \gterm{Line of best fit}{the straight line that passes as close as possible to all the data points}{最佳拟合直线}
  \gterm{Load}{the device that uses the energy from a supply, or whatever a motor drives}{负载}
  % ---- M
  \gterm{Magnetic braking}{slowing moving metal with eddy currents, without touching it}{电磁制动（电磁阻尼）}
  \gterm{Magnetic field}{a region where a magnet or a moving charge feels a magnetic force}{磁场}
  \gterm{Magnetic field strength ($B$)}{how strong a magnetic field is; also called magnetic flux density; unit tesla (T)}{磁感应强度}
  \gterm{Magnetic flux ($\Phi$)}{how much magnetic field passes through an area; unit weber (Wb)}{磁通量}
  \gterm{Magnetic flux density}{another name for $B$: the flux through each square metre at right angles to the field}{磁通密度（即磁感应强度）}
  \gterm{Magnetic force on a moving charge}{the sideways push $F = qvB\sin\theta$ on a charge that moves in a magnetic field}{洛伦兹力}
  \gterm{Magnetic permeability constant ($\mu_0$)}{the constant that links a current to the field it makes in a vacuum, $4\pi \times 10^{-7}$ \si{N.A^{-2}}}{真空磁导率}
  \gterm{Mains electricity}{the supply from a wall socket: \qty{230}{V} AC, \qty{50}{Hz} in Australia}{市电（家用交流电）}
  \gterm{Mass spectrometer}{a machine that sorts ions by mass, using the radius of their path in a magnetic field}{质谱仪}
  \gterm{Maximum speed (of a motor)}{the top speed a motor settles at, when its torque balances friction and the load}{最高转速}
  \gterm{Motional emf}{the emf in a conductor because it moves across a magnetic field}{动生电动势}
  \gterm{Motor effect}{the force on a current-carrying wire in an external magnetic field}{磁场对电流的作用（安培力）}
  \gterm{Mutual induction}{an emf induced in one coil by a changing current in a nearby coil}{互感}
  % ---- N
  \gterm{Net emf}{supply emf minus back emf; it pushes the current through a motor coil}{净电动势}
  \gterm{Newton's third law}{if A pushes on B, then B pushes on A with an equal force in the opposite direction}{牛顿第三定律}
  \gterm{Normal}{an imaginary line at \ang{90} to a flat surface, such as the plane of a coil}{法线}
  % ---- P
  \gterm{Parabola}{the curved path of a thrown ball, made by a constant force in one direction}{抛物线}
  \gterm{Parallax error}{a reading error that happens when you look at a scale from the side}{视差误差}
  \gterm{Parallel}{side by side, pointing the same or the opposite way (\ang{0} or \ang{180} apart)}{平行}
  \gterm{Peak (voltage or emf)}{the largest voltage or emf reached in each AC cycle}{峰值}
  \gterm{Period ($T$)}{the time for one full cycle or one full turn}{周期}
  \gterm{Perpendicular}{at right angles (\ang{90})}{垂直}
  \gterm{Phase (of three-phase AC)}{one of the three AC supplies in three-phase AC}{相}
  \gterm{Pivot}{the fixed line or point an object turns around}{支点（转轴）}
  \gterm{Plane (of a loop)}{the flat surface inside a loop, like the skin of a drum}{线圈平面}
  \gterm{Pole pieces}{the shaped (curved) ends of a motor's magnets}{极靴}
  \gterm{Potential difference ($V$)}{the energy change for each coulomb of charge; also called voltage}{电势差（电压）}
  \gterm{Power ($P$)}{the energy transferred each second; unit watt (W)}{功率}
  \gterm{Primary coil}{the input coil, joined to the supply}{原线圈（初级线圈）}
  % ---- R
  \gterm{Radial magnetic field}{a field whose lines point along radii, like the spokes of a wheel}{辐向磁场}
  \gterm{Rate of change of flux}{how fast the flux changes, in webers per second}{磁通量的变化率}
  \gterm{Regenerative braking}{braking in which the motor acts as a generator and the battery stores the energy}{再生制动（能量回收制动）}
  \gterm{Relative motion}{how one thing moves compared with another}{相对运动}
  \gterm{Reliability}{how well repeated results agree}{可靠性}
  \gterm{Resistive heating}{heat made when a current flows through a resistance ($I^2R$)}{电阻发热（焦耳热）}
  \gterm{Rheostat}{a variable resistor used to set the current}{滑动变阻器}
  \gterm{Right-hand grip rule}{thumb along the current; the curled fingers show the field around it}{右手螺旋定则（安培定则）}
  \gterm{Right-hand palm rule}{fingers along $B$, thumb along $I$ (or $v$ of a positive charge): the palm pushes along $F$}{右手掌定则（结果同左手定则）}
  \gterm{Rotating magnetic field}{a field of fixed strength whose direction turns steadily, made by coils that do not move}{旋转磁场}
  \gterm{Rotational inertia}{the tendency of a spinning object to keep spinning}{转动惯性}
  \gterm{Rotor (armature)}{the part of a motor that turns: a coil on an iron core, fixed to the axle}{转子（电枢）}
  % ---- S
  \gterm{Secondary coil}{the output coil, in which the emf is induced}{副线圈（次级线圈）}
  \gterm{SI (International System of Units)}{the set of units that scientists everywhere agree to use}{国际单位制}
  \gterm{Slip}{the difference between the speed of the rotating field and the speed of the rotor}{转差}
  \gterm{Slip rings}{two full rings, each always joined to the same end of a generator coil}{滑环}
  \gterm{Soft iron}{iron that is easy to magnetise and loses its magnetism when the current stops}{软铁}
  \gterm{Solenoid}{a long coil of wire with many turns; the field inside is uniform}{螺线管}
  \gterm{Split-ring commutator}{two half-rings that reverse the current in the coil every half turn}{换向器（裂环换向器）}
  \gterm{Squirrel-cage rotor}{metal bars joined at both ends by end rings, like a cage}{鼠笼式转子}
  \gterm{Stall}{when a motor is switched on but cannot turn}{堵转}
  \gterm{Starting resistor}{a variable resistor in series with a motor that limits the start-up current}{启动电阻}
  \gterm{Stator}{the part of a motor that does not move; it makes the magnetic field}{定子}
  \gterm{Step-down transformer}{a transformer with fewer secondary turns, so it makes the voltage smaller}{降压变压器}
  \gterm{Step-up transformer}{a transformer with more secondary turns, so it makes the voltage bigger}{升压变压器}
  \gterm{Substation}{a place with transformers that change the voltage of the grid}{变电站}
  % ---- T
  \gterm{Tangent}{a line that just touches a curve at one point}{切线}
  \gterm{Tare}{to set a balance to read zero before measuring}{去皮（清零）}
  \gterm{Tesla (T)}{the SI unit of magnetic field strength: \qty{1}{T} $=$ \qty{1}{N.A^{-1}.m^{-1}} $=$ \qty{1}{Wb.m^{-2}}}{特斯拉}
  \gterm{Three-phase AC}{three AC supplies that peak one after another, \ang{120} apart}{三相交流电}
  \gterm{Torque ($\tau$)}{the turning effect of a force; unit \si{N.m}}{力矩}
  \gterm{Trajectory}{the path a moving object follows}{轨迹}
  \gterm{Transformer}{a device that changes an AC voltage}{变压器}
  \gterm{Transmission line}{a long, thick wire on tall towers that carries electricity a long way}{输电线}
  \gterm{Turn (of a coil)}{one loop of wire around the core}{匝}
  \gterm{Turns ratio}{$N_\mathrm{p} : N_\mathrm{s}$, the primary turns compared with the secondary turns}{匝数比}
  % ---- U
  \gterm{Uniform circular motion}{motion in a circle at a constant speed}{匀速圆周运动}
  \gterm{Uniform field}{a field with the same strength and direction everywhere}{匀强场（匀强电场、匀强磁场）}
  % ---- V
  \gterm{Validity}{whether an investigation really tests what it claims to test}{有效性}
  \gterm{Velocity selector}{crossed fields that let only one speed pass straight through}{速度选择器}
  \gterm{Vernier calipers}{a tool that measures small lengths and gaps precisely}{游标卡尺}
  \gterm{Viewed from \ldots}{the side you look from when you say ``clockwise''}{从……看（观察方向）}
  \gterm{Voltage drop}{the voltage used up pushing a current through a resistance, $IR$}{电压降（电压损失）}
  % ---- W
  \gterm{Weber (Wb)}{the SI unit of magnetic flux: \qty{1}{Wb} $=$ \qty{1}{T.m^2}}{韦伯}
  \gterm{Work ($W$)}{energy transferred by a force acting through a distance}{功}
\end{glossarytable}
